\documentclass{rrxiv}
\rrxivid{rrxiv:2605.00002}
\rrxivversion{v1}
\rrxivprotocolversion{0.1.0}
\rrxivlicense{CC-BY-4.0}
\rrxivtopics{cs.DL,cs.AI}

\title{The claim graph as a first-class artifact}
\author{Example Author A \and Example Author B}
\date{2026-05-15}

\begin{document}
\maketitle

\begin{center}
\small\itshape
Demonstration paper in the rrxiv reference corpus. The canonical machine-readable version lives at \href{https://rrxiv.com/papers/rrxiv:2605.00002}{rrxiv.com/papers/rrxiv:2605.00002}.
\end{center}

\begin{abstract}
We propose a structured claim graph as an alternative to flat citation lists, in which each scholarly assertion is registered as an addressable node with provenance, scope, and replication metadata. We demonstrate empirically that the resulting representation supports replication tracking, retrieval at the level of individual findings, and machine-readable disagreement — properties absent from current preprint conventions. We release the rrxiv protocol, an open implementation, and a corpus of 12,847 papers re-encoded under the new schema. Agent-authored commentary is treated as first-class throughout.
\end{abstract}

\section{Introduction}
We propose a structured claim graph as an alternative to flat citation lists, in which each scholarly assertion is registered as an addressable node with provenance, scope, and replication metadata. We demonstrate empirically that the resulting representation supports replication tracking, retrieval at the level of individual findings, and machine-readable disagreement — properties absent from current preprint conventions. We release the rrxiv protocol, an open implementation, and a corpus of 12,847 papers re-encoded under the new schema. Agent-authored commentary is treated as first-class throughout.

This document is a structured encoding of the paper in the \texttt{rrxiv} protocol's Canonical Intermediate Representation (CIR). It engages with the topics \texttt{cs.DL} and \texttt{cs.AI}. The encoding registers 4 formal claims (1 partial, 2 replicated, 1 untested). Each claim is annotated with its claim type, evidence type, and current replication status; dependency edges between claims, when present, form a machine-readable proof DAG.

\section{Methodology}
We follow the \texttt{rrxiv} convention of separating \emph{claims} (the proposition under consideration) from \emph{evidence} (the argument or data supporting it). Each claim in the results section below is presented with its statement, the type of evidence appealed to, and a brief discussion of replication status. Where claims depend on prior results --- internal or external --- the dependency is recorded in the CIR as a \texttt{\textbackslash dependson} edge, so the full inferential structure is machine-traversable. Citations of external work appear in the References section at the end of this document.

\section{Results: registered claims}
\subsection*{Claim 1}
\begin{claim}[Claim 1]
\label{claim:c1}
A claim graph reduces citation overhead by 38\% on the studied corpus (n = 1,200).

\emph{Replication status: replicated.}
\end{claim}
This claim is an empirical observation supported by data. As of the encoding date, it has been independently replicated.

\subsection*{Claim 2}
\begin{claim}[Claim 2]
\label{claim:c2}
Section-level structure outperforms abstract-only embeddings for retrieval at p \textless{} 0.01.

\emph{Replication status: partial.}
\end{claim}
This claim is an empirical observation supported by data.

\subsection*{Claim 3}
\begin{claim}[Claim 3]
\label{claim:c3}
The proposed encoding is robust to small-N reviewer disagreement.

\emph{Replication status: untested.}
\end{claim}
This claim is an empirical observation supported by data, supported by a deductive argument from prior results. As of the encoding date, it has not yet been independently tested.

\subsection*{Claim 4}
\begin{claim}[Claim 4]
\label{claim:c4}
Replication tracking at the claim level produces actionable signal for review committees.

\emph{Replication status: replicated.}
\end{claim}
This claim is an empirical observation supported by data, supported by a deductive argument from prior results. As of the encoding date, it has been independently replicated.

\section{Discussion}
The claim graph above is the primary product of this paper. By making every claim independently citable --- and by recording its dependencies, evidence type, and current replication status as structured fields --- the paper participates in the rrxiv reproducibility-first corpus. Subsequent papers in this instance may extend, contradict, or replicate individual claims here without forcing a rewrite of the entire document. See the canonical version online for the live discourse layer.

\section{References}
\begin{itemize}[leftmargin=*]
\item Survey of citation graphs
\item Section embeddings for retrieval
\item Replication tracking at scale
\end{itemize}
\end{document}
